\chapter{公式速查表}
\small
\section{固体物理}
\begin{longtable}{p{0.42\textwidth}p{0.5\textwidth}}\toprule
公式&物理意义\\\midrule\endhead
$\vb a_i\cdot\vb b_j=2\pi\delta_{ij}$&定义倒格基矢，保证晶格平移相位为 $2\pi$ 整数倍。\\
$\vb k'-\vb k=\vb G\Leftrightarrow2d\sin\theta=n\lambda$&Laue 与 Bragg 衍射条件。\\
$S_G=\sum_jf_j e^{-iG\cdot r_j}$&基元内散射振幅及系统消光。\\
$\omega=2\sqrt{K/M}|\sin(ka/2)|$&一维单原子链声子色散。\\
$g_{ph}(\omega)\propto\omega^2,\ C_V^{ph}\propto T^3$&三维低频声学模决定低温热容。\\
$g(E)=\frac{V}{2\pi^2}(2m/\hbar^2)^{3/2}\sqrt E$&含自旋的三维自由电子态密度。\\
$E_F=\frac{\hbar^2}{2m}(3\pi^2n)^{2/3}$&零温自由电子最高占据能。\\
$C_V^{el}=\frac{\pi^2}{3}\kb^2g(E_F)T$&费米面附近电子给线性热容。\\
$\psi_{nk}=e^{ik\cdot r}u_{nk},\ \Delta E=2|V_G|$&Bloch 形式与区边界开隙。\\
$v=\hbar^{-1}\nabla_kE,\ \hbar\dot k=-eE$&波包的能带半经典动力学。\\
$np=n_i^2$&平衡非简并半导体质量作用定律。\\
$\chi=C/(T-T_C)$&局域磁矩平均场响应。\\
$\nabla^2B=B/\lambda_L^2$&超导体内磁场指数衰减。\\\bottomrule
\end{longtable}

\section{理论力学}
\begin{longtable}{p{0.42\textwidth}p{0.5\textwidth}}\toprule
公式&物理意义\\\midrule\endhead
$\delta S=0,\ \frac d{dt}\partial_{\dot q_i}L-\partial_{q_i}L=0$&作用量驻值给运动方程。\\
$p_i=\partial_{\dot q_i}L,\ H=\sum p_i\dot q_i-L$&Legendre 变换到相空间。\\
$\dot q_i=\partial_{p_i}H,\ \dot p_i=-\partial_{q_i}H$&Hamilton 正则方程。\\
$\dot f=\{f,H\}+\partial_tf$&Poisson 括号生成时间演化。\\
$\det(K-\omega^2M)=0$&多自由度小振动的正规模条件。\\
$\dot{\vb L}_{body}+\vb\omega\times\vb L=\vb N$&随体坐标中的刚体 Euler 方程。\\\bottomrule
\end{longtable}

\section{电动力学}
\begin{longtable}{p{0.42\textwidth}p{0.5\textwidth}}\toprule
公式&物理意义\\\midrule\endhead
$\nabla\cdot E=\rho/\varepsilon_0,\ \nabla\cdot B=0$&场的源：电荷存在、磁单极缺席。\\
$\nabla\times E=-\partial_tB$&变化磁场产生涡旋电场。\\
$\nabla\times B=\mu_0J+\mu_0\varepsilon_0\partial_tE$&传导与位移电流产生磁场。\\
$\nabla^2\phi=-\rho/\varepsilon_0$&静电势的 Poisson 方程。\\
$\nabla^2E-c^{-2}\partial_t^2E=0$&真空电磁场以光速传播。\\
$S=E\times H$&单位面积单位时间的电磁能流。\\
$A'=A+\nabla\chi,\ \phi'=\phi-\partial_t\chi$&势有规范冗余而物理场不变。\\
$\beta^2=\mu\varepsilon\omega^2-k_c^2$&波导传播常数与截止本征值。\\
$Z_{\rm TE}=\omega\mu/\beta,\ Z_{\rm TM}=\beta/(\omega\varepsilon)$&TE、TM 模式的横向场阻抗。\\
$\partial_\mu F^{\mu\nu}=4\pi J^\nu$&$c=1$ 的 Gaussian 几何制非齐次协变 Maxwell 方程。\\
$E^2-B^2,\ \vb E\cdot\vb B$&Gaussian 几何制的两个 Lorentz 场不变量。\\\bottomrule
\end{longtable}

\section{热力学与统计物理}
\begin{longtable}{p{0.42\textwidth}p{0.5\textwidth}}\toprule
公式&物理意义\\\midrule\endhead
$dU=TdS-PdV+\mu dN$&简单系统的热力学基本关系。\\
$F=U-TS,\ G=U-TS+PV,\ \Omega=U-TS-\mu N$&适配不同控制变量的热力学势。\\
$P_i=e^{-\beta E_i}/Z$&系统与热库平衡的正则分布。\\
$F=-\kb T\ln Z,\ U=-\partial_\beta\ln Z$&配分函数生成热力学量。\\
$PV=N\kb T$&经典理想气体状态方程。\\
$\bar n=[e^{\beta(\epsilon-\mu)}\mp1]^{-1}$&Bose（减）与 Fermi（加）平均占据。\\
$S=\kb\ln\Omega$&微正则态数定义孤立系统熵。\\
$\rho_{\rm can}=e^{-\beta H}/Z$&最大 von Neumann 熵给正则热密度矩阵。\\
$\rho_{\rm gc}=e^{-\beta(H-\mu N)}/\Xi$&巨正则密度矩阵同时约束平均能量和粒子数。\\
$S=-\kb\operatorname{Tr}(\rho\ln\rho)$&von Neumann 熵统一纯态、混合态与系综。\\
$\Xi=\operatorname{Tr}e^{-\beta(H-\mu N)}$&巨正则系综允许交换能量与粒子。\\
$\langle(\Delta E)^2\rangle=\kb T^2C_V$&正则能量涨落与热容的联系。\\
$N_0/N=1-(T/T_c)^{3/2}$&三维均匀理想 Bose 气凝聚分数。\\
$F=F_0+a(T-T_c)m^2+bm^4$&Landau 平均场描述连续相变和对称破缺。\\\bottomrule
\end{longtable}

\section{量子力学}
\begin{longtable}{p{0.42\textwidth}p{0.5\textwidth}}\toprule
公式&物理意义\\\midrule\endhead
$i\hbar\partial_t\psi=H\psi$&封闭量子态的幺正时间演化。\\
$\partial_t\rho+\nabla\cdot j=0$&概率局域守恒。\\
$E_n=\hbar\omega(n+\frac12)$&谐振子等间隔能级与零点能。\\
$[J_i,J_j]=i\hbar\epsilon_{ijk}J_k$&旋转生成元的代数结构。\\
$P_{12}\Psi_\pm=\pm\Psi_\pm$&全同粒子的交换对称或反对称扇区。\\
$\Psi_\pm=(\phi_a(1)\phi_b(2)\pm\phi_b(1)\phi_a(2))/\sqrt{2(1\pm|S|^2)}$&非正交轨道的两粒子对称化。\\
$\langle V\rangle_\pm=J_{ab}\pm K_{ab}$&直接相互作用与交换干涉共同决定两粒子能量。\\
$\{c_i^\dagger,c_i^\dagger\}=0\Rightarrow(c_i^\dagger)^2=0$&费米反对易关系编码 Pauli 不相容。\\
$E_n^{H}=-13.6\,{\rm eV}/n^2$&非相对论氢原子束缚能级。\\
$E_n^{(1)}=\langle n|H'|n\rangle$&非简并微扰的一阶能移。\\
$E_n^{(2)}=\sum_{m\ne n}|\langle m|H'|n\rangle|^2/(E_n-E_m)$&虚跃迁导致的二阶能移。\\
$E[\psi]\ge E_0$&任意试探态能量给基态上界。\\
$\int_{x_1}^{x_2}p\,dx=(n+\frac12)\pi\hbar$&两转向点 WKB 半经典量子化。\\
$\Delta A\Delta B\ge\frac12|\langle[A,B]\rangle|$&不对易导致的内禀涨落下界。\\
$\Delta E_Z=\mu_Bg_Jm_jB$&弱场反常 Zeeman 能级劈裂。\\
$\Delta E_{\rm Stark}^{(2)}=-\alpha\mathcal E^2/2$&非简并定态的二次 Stark 位移。\\
$d\sigma/d\Omega=|f(\theta,\phi)|^2$&散射振幅决定微分截面。\\
$f_B(\vb q)=-\frac{m}{2\pi\hbar^2}\int e^{-i\vb q\cdot\vb r}V(\vb r)d^3r$&一阶 Born 振幅是势的 Fourier 分量。\\
$f=\frac1k\sum_l(2l+1)e^{i\delta_l}\sin\delta_lP_l$&中心势散射的分波展开。\\
$\sigma_{\rm tot}=4\pi\,\operatorname{Im}f(0)/k$&光学定理表达概率流守恒。\\\bottomrule
\end{longtable}
\normalsize

\section{最后一分钟检查}
每个回答至少包含：定义；一条关键公式；适用条件；一个物理图像。若公式记不全，先说明量纲、极限和对称性，通常比猜系数更可靠。
